\documentclass{amsart}
\usepackage{mathtools}
\usepackage{glossaries-extra}
\usepackage[top=2cm, bottom=2cm, left=2cm, right=2cm]{geometry}
\usepackage{xcolor}

\renewcommand{\d}{\mathrm d}

\title{Sharp Large Deviations for the Elephant Random Walk}
\date{November 18th, 2026}
\author{Juan Esaul González Rangel}

\newacronym{erw}{ERW}{Elephant Random Walk}
\newacronym{ldp}{LDP}{Large Deviation Principle}

\usepackage{biblatex}
\bibliography{references}
\usepackage{cleveref}

\begin{document}

\maketitle

Let $S=(S_n, n \ge0)$ be an \gls{erw} with memory parameter $p$ and first step parameter $q$, and define \[ \mu \coloneqq \lim_{n\to\infty} \mathbb E \left[ S_n/n \right] < \infty. \] 
For any $j = n + 2 k$ with $k\in \mathbb Z$ and every $t>0$, we have
\begin{align*}
    \mathbb P ( S_n = j ) &= \mathbb P \left(\frac{S_n - n}2 = k \right) = \frac{e^{-tk}}{2\pi}\int_{-\pi}^\pi e^{-iuk} \mathbb E[ e^{(t+iu)\frac{s_n-n}{2}} ] \d u \\
    &= \frac{1}{2\pi} \int_{-\pi}^\pi e^{-(t+iu)\left( k - \frac n2 \right)} m_n\left(\frac{t+iu}{2}\right) \mathrm d u\intertext{taking $v=(u +it)/2$ and $k = (j-n)/2$}
    &= \frac1\pi \int_{-\frac\pi2}^{\frac\pi2} e^{-(t+iv)j}m_n(t+iv) \d v.
\end{align*}
So that, if $x \in (\mu,1)$, we have 
\begin{align}
    \mathbb P( S_n \ge nx ) &= \sum_{j=\lceil nx \rceil}^n \frac{e^{-tj}}{\pi}\int_{-\frac\pi2}^{\frac\pi2} e^{-ivj} m_n(t+iv) \mathrm d v = \sum_{j= \lceil nx \rceil}^\infty \frac{e^{-tj}}{\pi}\int_{-\frac\pi2}^{\frac\pi2} e^{-ivj} m_n(t+iv) \mathrm d v \notag \\
    &= \frac1{\pi} \int_{-\frac\pi2}^{\frac\pi2} \frac{ e^{-(t+iv)\lceil nx \rceil} }{1 - e^{-(t+iv)}}m_n(t+iv) \mathrm d v \notag \\
    &= \frac{ e^{ -t\lceil nx \rceil } \phi^n(t) }{\pi} \int_{-\frac\pi2}^{\frac\pi2} \frac{ e^{-i v \lceil nx\rceil} }{ 1- e^{-(t+iv)} } \frac{m_n(t+iv)}{\phi^n(t+iv)} \left(\frac{\phi(t+iv)}{\phi(t)}\right)^n \mathrm d v \notag \\
    &= \frac{ e^{-n(xt-\Lambda(t))- \left\{ nx\right\}t } }{\pi} \left( \int_{-\frac\pi2}^{\frac\pi2} \frac{e^{-n\varphi(v)}}{1-e^{-(t+iv)}}\mathrm d v +  \int_{-\frac\pi2}^{\frac\pi2} \frac{e^{-n\varphi(v)}}{1-e^{-(t+iv)}}r_n(t+iv)\mathrm d v\right), \label{eq:integral}
\end{align}
where \[\Lambda(z) \coloneqq \log \phi(z), \qquad \varphi(v) \coloneqq iv \frac{\lceil nx\rceil}{n} - \Lambda(t+iv) + \Lambda(t), \] and \[ r_n(t+iv) \coloneqq \frac{m_n(t+iv)}{\phi^n(t+iv)} - 1. \]

From now on, we must take $t = t_x \coloneqq \arg\max_{t>0} \left( xt - \Lambda(t) \right)$, which is the unique solution to \[ x = \Lambda'(t_x), \] and it exists since $\Lambda$ is strictly convex. The reason for such a choice comes from the fact that the exponential rate in \cref{eq:integral} is indeed $\Lambda^*(x)\coloneqq \sup_{t>0} (xt - \Lambda(t))$, by the known \gls{ldp} for the \gls{erw} (see eg. \cite{franchini2023large}).

For the first integral in \cref{eq:integral}, we have, for any $0<\epsilon<\pi/2$

\begin{align*}
    \int_{-\frac\pi2}^{\frac\pi2} \frac{e^{-n\varphi(v)}}{1-e^{-(t_x+iv)}}\mathrm d v &= \int_{-\epsilon}^\epsilon \frac{e^{-n\varphi(v)}}{1-e^{-(t_x+iv)}}\mathrm d v + \int_{[-\epsilon,\epsilon]^c} \frac{e^{-n\varphi(v)}}{1-e^{-(t_x+iv)}}\mathrm d v \eqcolon I_1 + I_2.
\end{align*}
We can, thus, choose $\epsilon$ small enough. To deal with the term $I_1$
\begin{align*}
    I_1 &= \int_{-\epsilon}^\epsilon \frac{e^{-n ( ivx - \Lambda(t_x+iv ) + \Lambda(t_x) )}}{1-e^{-(t_x + iv)}}\mathrm d v = \frac{1}{\sqrt{n}}\int_{-\epsilon\sqrt{n}}^{\epsilon\sqrt{n}} \frac{e^{-n ( i \frac{u}{\sqrt{n}}x - \Lambda(t_x+i\frac{u}{\sqrt{n}} ) + \Lambda(t_x) )}}{1-e^{-(t_x +i\frac{v}{\sqrt{n}})}}\mathrm d u. \intertext{A Taylor expansion of $\Lambda$ around $t_x$ yields}
    I_1 &= \frac{1}{\sqrt{n}}\int_{-\epsilon\sqrt{n}}^{\epsilon\sqrt{n}} \frac{e^{-n \left(\frac{u^2}{2n}\Lambda''(t_x) + \frac1n\eta\left(\frac{u}{\sqrt{n}}\right) \right)}}{1-e^{-(t_x + i\frac{u}{\sqrt{n}})}}  \mathrm d u = \frac{1}{\sqrt{n}} \int_{-\epsilon\sqrt{n}}^{\epsilon\sqrt{n}} \frac{e^{-\frac{u^2}{2}\Lambda''(t_x) + \eta\left(\frac u{\sqrt{n}}\right) }}{1-e^{-(t_x + i \frac{u}{\sqrt{n}})}} \mathrm d u, \intertext{where, by hypothesis, $\eta\left( \frac{u}{\sqrt{n}} \right)\to 0$ as $n\to\infty$.}
\end{align*}
Notice that \[  \left\lvert\int_{-\epsilon\sqrt{n}}^{\epsilon\sqrt{n}} \frac{e^{-\frac{u^2}{2}\Lambda''(t_x) + \eta\left(\frac u{\sqrt{n}}\right) }}{1-e^{-(t_x + i \frac{u}{\sqrt{n}})}}\mathrm d u \right\rvert \le \int_{-\infty}^{\infty} \frac{e^{-\frac{u^2}{2}\Lambda^{''}(t_x) + \xi}}{\sqrt{2e^{-t_x}\cosh t_x}} \d u < \infty,\]
where \[\xi \coloneqq \sup\left\{ \lvert\eta(u/\sqrt n)\rvert :   n\in \mathbb N, u \in [-\epsilon\sqrt{n}, \epsilon\sqrt{n}]  \right\} = \sup\left\{ \lvert\eta(u)\rvert : u \in [ -\epsilon, \epsilon] \right\} < \infty, \] which is finite since $\eta$ is the difference between two continuous functions in $[-\pi,\pi].$

The Dominated Convergence Theorem then yields \[ \lim_{n\to\infty} \sqrt{n} I_1 = \frac{1}{1-e^{-t_x}}\sqrt{ \frac{2\pi}{\Lambda^{''}(t_x)} }.  \]

The term $I_2$ can be bounded as follows:
\begin{align*}
    |I_2| &\le \int_{[-\epsilon,\epsilon]^c} \frac{\left| e^{-n (\Lambda(t_x) - \Lambda(t_x+iv))} \right|}{\sqrt{2e^{-t_x} (\cosh t_x -\cos v) }} \mathrm d v \le \int_{[-\epsilon,\epsilon]^c} \frac{e^{-n (\Lambda(t_x) - \Re(\Lambda(t_x+iv))} }{\sqrt{2e^{-t_x}\cosh(t_x)}} \mathrm d v \intertext{Take $\delta(\epsilon) \coloneqq \inf\{ \Lambda(t_x) - \Re(\Lambda(t_x+iv)) : v \in [-\epsilon, \epsilon]^c \} > 0$, so }
    \lvert I_2 \rvert &\le \frac{(\pi-2\epsilon) e^{-n \delta(\epsilon)}}{\sqrt{2e^{-t_x}\cosh(t_x)}}, \intertext{and then}
    \sqrt{n}I_2 &\xrightarrow[n\to\infty]{\phantom{yeii}} 0.
\end{align*}
Together with the convergence for $I_1$, this implies that as $n\to \infty$ \[ \int_{-\frac\pi2}^\frac\pi2 \frac{e^{-n\varphi(u)} }{1-e^{-(t_x + iv)}}\d v = \sqrt{\frac{2\pi}{n\Lambda^{''}(t_x)}}\frac{\left(1 + o(1)\right)}{1 - e^{-t_x}}.\]


Now, we need to deal with the second integral in \cref{eq:integral}. 

\begin{align*}
    \int_{-\frac\pi2}^{\frac\pi2} \frac{e^{-n\varphi(v)}}{1-e^{-(t+iv)}}r_n(t+iv)\mathrm d v
\end{align*}

\printbibliography

\section*{Computations}



\end{document}